\chapter{Self-consistency and numerical reliability}
\label{narr:variation}

A self-consistency equation characterizes a candidate state in a specified
approximation. Solving it accurately is necessary for that approximation, but
does not by itself identify the physical ground state. This chapter separates
three questions: whether the state is admissible, whether it is stationary, and
whether it is energetically or dynamically stable under the variations allowed.

% EDITORIAL QUESTION: What is the geometric stationarity equation for a pure Gaussian trial state?
\section{Variation on the pure covariance manifold}
Let $E_G(\Omega)$ be the expectation of a declared auxiliary Hamiltonian in a
pure Gaussian state, evaluated with the convention of Chapter~\ref{narr:gaussian}.
Suppose smooth constraints $g_a(\Omega)=0$ are imposed and define
$\mathcal L=E_G+\sum_a\lambda_ag_a$. Define the real skew gradient $A_{\mathcal L}$
by the differential convention
\begin{equation}
 \delta\mathcal L=\tfrac14\Tr(A_{\mathcal L}^T\delta\Omega).
 \label{eq:narr:gradient}
\end{equation}
An allowed unconstrained orbit variation is
$\Omega(t)=e^{tX}\Omega e^{-tX}$ with $X^T=-X$, so its tangent is $[X,\Omega]$.
Stationarity of the Lagrangian gives
\begin{equation}
 [A_{\mathcal L},\Omega]=0,\qquad g_a(\Omega)=0.
 \label{eq:narr:stationarity}
\end{equation}
Appendix~\ref{app:finite_foundations} derives the tangent description and this
commutator condition. The multiplier formulation presupposes the usual
constraint regularity needed to represent constrained stationarity this way.
A singular constraint set may require a direct tangent-cone analysis instead.

Equation~\eqref{eq:narr:stationarity} is necessary, not a minimum certificate.
Commuting matrices may describe maxima or saddles, and a spectral occupation
choice is additional information. Even a local minimum within Gaussian states
need not be a global Gaussian minimum or a physical spin minimum. In particular,
weak local constraints do not automatically implement Eq.~\eqref{eq:narr:projection}.

% EDITORIAL QUESTION: What does a converged iteration establish?
\section{Residuals, competing branches, and stability}
An iteration chooses a map $\Phi$ and searches for a zero of
$R(\Omega)=\Phi(\Omega)-\Omega$. Its step count measures the chosen algorithm,
not physical relaxation time. Damping or mixing changes iteration stability
without changing the intended stationary equations when implemented consistently.
A small update can also result from a small damping factor while the true
residual remains large.

A useful verification records the equation residual, constraint residual,
antisymmetry/purity defects, energy, geometry, and stopping rule separately.
The relevant norm and its normalization matter as the system size grows.
Near a small spectral gap, perturbations can rotate an occupied subspace much
more than the eigenvalues suggest. Precision and tolerances should therefore
be tied to conditioning and observed numerical error rather than chosen once
for every calculation.

Multiple seeds and continuation from different directions can find competing
stationary branches. Hysteresis of a solver is evidence about branch tracking;
a thermodynamic first-order transition additionally requires a comparison of
the appropriate free energies and controlled size dependence. A branch losing
iteration convergence is not automatically the physical spinodal. A second
variation on admissible directions is the separate local stability question.
The historical solver chapters retain detailed experiments, but their numeric
thresholds are not recertified by this structural explanation.

% EDITORIAL QUESTION: When do upper and lower energy estimates bound the same problem?
\section{A matched bracketing problem}
A physical projected trial state gives an upper bound through
Eq.~\eqref{eq:narr:variational}. A lower bound can instead arise from relaxing
the physical moment constraints. For chosen operators $O_a$, every physical
state has a positive moment matrix
\begin{equation}
 M_{ab}=\Tr(\rho_s O_a^\dagger O_b),\qquad M\succeq0.
 \label{eq:narr:moment}
\end{equation}
If a feasible set contains all physical moment data and the objective equals
the physical energy on those data, its \emph{infimum} cannot exceed $E_0$.
This inclusion is the mathematical reason for a relaxation lower bound.

A numerical feasible point in the relaxed primal problem is an upper bound
on that relaxed infimum; it need not be a lower bound on the physical energy.
A certified dual-feasible value, with numerical errors controlled, is one route
to an actual reported lower bound. Thus a table headed ``SDP lower bound'' needs
both the relaxation contract and a valid certificate. Bracketing is meaningful
only when upper and lower constructions refer to the same Hamiltonian, physical
space, geometry, units, and sector restrictions.
